\paragraph{A product CM point on each indefinite marked row.}
For the fields of Proposition~\ref{host:art} there is another
choice of evaluating point. Write \(k=\mathbb Q(\sqrt{-d_0})\),
\(F_0=E^+\), \(O=O_E\), and \(\Phi\) for the embeddings inducing
the chosen embedding of \(k\). For \(\tau\in\Phi\) write \(p,q\)
for the Lie multiplicities at \(\tau,\bar\tau\), respectively.
Write \(p_0\) for the number of positive Hermitian row directions.
The selected period map conjugates the first \(p_0=q\) coordinates
and is \(\tau\)-linear on the last \(r-p_0=p\) coordinates.
Keep the period convention \(\Pi=\mathfrak a L\) of
Lemma~\ref{bfam:minimum}, so \(L\) is Hermitian self-dual.
\src{src/Formalism/LocalToGlobal/ProjectiveLattices/PetersonNormSubspace.lean}{67}{The conjugated positive block in the period map};
\src{src/Formalism/LocalToGlobal/ProjectiveLattices/PetersonNormSubspace.lean}{75}{The positive row index on a good chart};
\src{src/ReciprocityTheory/ReductiveGroups/ArithmeticGeometry/BirationalGeometry/ToroidalCompactification/FlagBasis.lean}{76}{The Lie multiplicities of the two embeddings}.

\begin{lemma}\label{cmeig:row}
If \(p,q\ge1\), every marked row in this family has, up to rational
isometry, underlying lattice
\[
        J e_1\ \perp\ O e_2\ \perp\cdots\perp\ O e_r,\qquad
        J\bar J=O,\qquad
        h(e_i,e_i)=
        \begin{cases}1&i\le p_0,\\-1&i>p_0.\end{cases}
\]
The transported mark need not split into marks of the summands.
\end{lemma}
\begin{proof}
The two imaginary quadratic subfields of discriminants
\(-d_0,-q_0\) in Proposition~\ref{host:art} have coprime
discriminants (here \(q_0\) is the auxiliary prime). At each
rational prime at least one of them is unramified. Its inertia
restriction is trivial, whereas complex conjugation acts on it
nontrivially. Thus \(E/F_0\) is split or unramified at every finite
place, since relative inertia is the intersection with
\(\operatorname{Gal}(E/F_0)\).

At a split place a basis of the first factor and its pairing-dual
in the second identify a self-dual lattice with the standard one.
At an unramified quadratic place its reduction is a nondegenerate
Hermitian space over the quadratic residue extension. There is a
vector of nonzero norm: otherwise polarization of \(h(x+ay,x+ay)\)
and the perfect residue trace pairing would give \(h=0\).
Lifting gives a unit norm, which can be scaled to either sign,
since the norm maps unramified units onto units (use the residue
norm and lift successively with the residue trace). Its integral
orthogonal complement is again self-dual. This gives the local
orthogonal basis by induction, including in residue characteristic
two. At infinity the signatures agree with the displayed form.
For clarity, the rational isometry follows already from
Hasse--Minkowski: the underlying quadratic form \(h(x,x)\)
represents the prescribed first sign at every place, hence globally.
For the latter implication apply isotropy to \(h(x,x)-a t^2\);
if \(t=0\) gives a nonzero isotropic vector, the original quadratic
space contains a hyperbolic plane and represents \(a\).
Split off the Hermitian line of norm \(a\); by Witt cancellation its
complement is locally isometric to the remaining diagonal form.
Induction now identifies the two rational Hermitian spaces.

Put \(L_0=\bigoplus O e_i\), \(G=U(h)\). Choose
\(g\in G(\mathbb A_{F_0,f})\) carrying \(\widehat L_0\) onto
\(\widehat L\) using the local isometries, choosing the identity
where the two lattices coincide. Write \(t=\det g\) and
\(d(t)=\operatorname{diag}(t,1,\dots,1)\). The determinant-one
groups have strong approximation away from the infinite places:
\(\mathrm{SU}(h)\) is simply connected, absolutely almost simple
(\(r\ge2\)), and noncompact at a real place by \(p,q\ge1\).
This is \href{https://www.mathnet.ru/links/02c4d3528af06a248862fb1aa3ba53f9/im2220_eng.pdf}
{Platonov, Theorem A}. Thus the open set
\(g K_0 d(t)^{-1}\cap\mathrm{SU}(h)(\mathbb A_{F_0,f})\)
contains a rational element \(\gamma\): it contains
\(g d(t)^{-1}\), and \(K_0\) is the open stabilizer of
\(\widehat L_0\). Then \(\gamma^{-1}L\) has completion
\(d(t)\widehat L_0\). Its first ideal is
\(J=E\cap t\widehat O\); \(t\bar t=1\) gives \(J\bar J=O\).
Transport the entire given finite marking by \(\gamma^{-1}\).
The resulting polarized marked row is isomorphic to the one
started with, as in Lemma~\ref{munf:marks}. No approximation to
the identity of the marking, or change of its level, is needed.
\end{proof}

\begin{proposition}[The degree on the product point]\label{cmeig:degree}
On each classical marked period component of
Lemma~\ref{cmeig:row} there is a point over a finite extension of
the coefficient field with everywhere good abelian reduction. For
every \(\tau\in\Phi\) its two integral Lie determinants in the
polarization metric of the cover satisfy
\[
 \widehat{\deg}T_+-\widehat{\deg}T_-=
   \frac{p-q}{2}\left(\gamma_{\rm Eul}+
     \log\frac{2\pi}{{\rm rd}(E)}
       -\frac{L'(1,\chi)}{L(1,\chi)}\right).
\]
All finite places, including those ramified in \(E\), occur with
their ordinary integral eigensublattices.
\end{proposition}
\begin{proof}
In the representative of Lemma~\ref{cmeig:row} take the flag of
the positive and negative coordinate axes. The period lattice
is the orthogonal direct sum of the fractional ideals
\(I_1=\mathfrak a J,\ I_i=\mathfrak a\ (i>1)\).
The last \(p\) factors have CM type \(\Phi\), and the first \(q\)
have its conjugate, by the period map above.
Let \(C\) be the elliptic curve with complex lattice \(O_k\),
with the indicated action (conjugate the action for the opposite
type). It is defined with its endomorphisms over a number field
and acquires good reduction everywhere after finite extension:
its CM \(j\)-value is an algebraic integer, hence gives potential
good reduction at the finitely many places of bad reduction.
The Serre construction \(I_i\otimes_{O_k} C\) exists integrally
since \(I_i\) is projective over \(O_k\): an idempotent on a
finite direct power of \(C\) realizes it as a direct factor.
Its periods at the chosen embedding are exactly the images of
\(I_i\) at the extensions of that \(k\)-embedding. The alternating
trace polarization on each factor is induced from a positive
Hermitian form \(\operatorname{Tr}_{E/k}(\beta u\bar v)\).
Indeed divide its purely imaginary trace coefficient (with the
coordinate sign) by the coefficient of the standard principal
elliptic polarization, a signed generator of
\(\mathfrak D_{k/\mathbb Q}^{-1}\); the ratio \(\beta\in E^+\)
is totally positive by the Riemann positivity condition.
Perfection is exactly
\[
       (\beta)\,I_i\bar I_i\,\mathfrak D_{E/k}=O .
\]
This also follows by taking the trace-dual ideal. Hence the
polarized product has exactly the required period lattice and
polarization. For the polarized tensor construction see
\href{https://arxiv.org/abs/1507.07607}{Amir-Khosravi, \S2}.
The transported finite torsion marking is defined after another
finite extension, irrespective of how it mixes factors. It is
needed on the characteristic-zero point; the underlying product
extends over the full ring of integers. In a classical PEL moduli
chart this constructs the point of the prescribed component,
including its canonical rigidified graph line. An additional
finite étale marking cover is handled by lifting the point to
the given component of the cover, whose map to the underlying
connected component is finite étale and surjective.

Here is the integral and metric calculation for a single factor
with ideal \(I\) and tensor coefficient \(\beta\). Enlarge the
number field \(K\) of definition to split \(E\), and put
\(\mathfrak D=\mathfrak D_{E/k}\). In
\(I\otimes_{O_k}O_K\), the intersection with the coordinate line
at an active embedding \(\nu\) is \(\nu(I\mathfrak D)O_K\).
In fact test membership in the lattice by its dual. The trace
dual of \(I\) over \(O_k\) is \(I^{-1}\mathfrak D^{-1}\);
this remains the dual after flat base change. On the split
generic algebra the trace is the sum of coordinates. The vector
supported at \(\nu\) with entry \(z\) thus pairs integrally
exactly when \(z\nu(I^{-1}\mathfrak D^{-1})\subset O_K\), as
asserted. Tensoring with the invertible \(O_K\)-module
\(\operatorname{Lie} C\) gives the actual Lie eigenline of the
abelian scheme, since taking Lie commutes with the idempotent
defining the tensor construction.

On that coordinate its squared metric at \(\sigma:K\to\mathbb C\)
is the squared elliptic Lie metric multiplied by
\(\sigma\nu(\beta)\): diagonalizing the trace Hermitian form in
the embedding coordinates gives coefficients \(\sigma\nu(\beta)\);
the polarization convention multiplies both Hermitian forms by
the same \(\pi\). For \(e=[E:\mathbb Q]\), normalized ideal and
metric degrees therefore give
\[
\begin{split}
 2\widehat{\deg}T_\nu(I)
 &=2\widehat{\deg}\operatorname{Lie}C
   -\frac{2\log N(I\mathfrak D)+\log|N_{E/\mathbb Q}\beta|}{e}\\
 &=2\widehat{\deg}\operatorname{Lie}C-\frac{\log N\mathfrak D}{e}.
\end{split}
\]
The second equality uses perfection, so even the fractional ideal
on the first coordinate disappears. The norm is the absolute
ideal norm and \(N\mathfrak D=|\Delta_E|/d_0^{e/2}\).

For \(C(\mathbb C)=\mathbb C/\Pi_C\) with Euclidean period area
\(a_C\), principal polarization gives the Hermitian coefficient
\(1/a_C\). The tangent metric is \(\pi/a_C\), so its dual on
\(dz\) is \(a_C/\pi=(i/2\pi)\int dz\wedge d\bar z\).
Thus \(-\widehat{\deg}\operatorname{Lie}C\) is exactly the
stable elliptic Faltings height with the normalization of
\href{https://people.math.wisc.edu/~tonghaiyang/KRY2.pdf}
{Kudla--Rapoport--Yang, (10.14)--(10.15), Proposition 10.10}.
Their quadratic CM formula gives
\[
 2\widehat{\deg}\operatorname{Lie}C
 =\gamma_{\rm Eul}+\log(2\pi/\sqrt{d_0})
                  - L'(1,\chi)/L(1,\chi).
\]
This applies to either elliptic orientation. Substitution gives
the same \(2\widehat{\deg}T_\nu(I)=A_\chi\) on every active
coordinate; orthogonality and the integral direct product give
the two determinants by addition over \(p,q\) coordinates.
\end{proof}

In particular no comparison between a mixed isogeny image and
its saturation is needed to evaluate this point. If the vertical
orders on a selected component agree at all tests and its
archimedean norm is constant, the degree is that of this point
by the normalized base-change calculation preceding
\eqref{ctwist:degree}.

\paragraph{Vertical orders at hyperspecial places.}
There is a direct good-reduction comparison for the ordinary Lie
frames of the classical PEL problem. It also distinguishes a prime
in the level from a prime ramified in the endomorphism algebra.

\begin{proposition}\label{cmeig:hyperspecial}
Suppose the characteristic-zero cover on the chosen component
is a classical finite marking level for its principally polarized
\(O_E\)-abelian scheme, and its finite tests use the integral Lie
eigensublattices of connected semiabelian extensions of that same
scheme. Let \(w\mid\ell\) in the splitting coefficient field \(F\),
with \(\ell\nmid 2\Delta_E\). Each rational horizontal determinant
generator has the same normalized order at all semiabelian DVR
tests on the component above \(w\). If its marked level itself is
hyperspecial at \(\ell\), this common order belongs to \(\mathbb Z\).
\end{proposition}
\begin{proof}
First take a neat prime-to-\(\ell\) marking with full hyperspecial
level at \(\ell\). The algebra is unramified at \(\ell\), its
involution is of unitary type and the principal alternation is
perfect. Thus this is a good reduction PEL problem. With smooth
projective cone data it has a proper smooth toroidal model with
universal semiabelian scheme by
\href{https://www.kwlan.org/articles/cpt-PEL-type-thesis-revision-dense.pdf}
{Lan, Theorems 6.4.1.1 and 7.3.3.4}.
Base change to the completed coefficient DVR. The geometric
components of a proper smooth family form a finite étale scheme;
the selected geometrically connected generic component defined
over the coefficient field extends to an open and closed component
with geometrically connected special fibre. That fibre is smooth,
so is reduced and irreducible. The split embedding idempotents of
\(O_E\otimes_{\mathbb Z}O_{F,w}\) are integral. The indicated Lie
determinant ratio is therefore a line bundle on the model.
Its characteristic-zero horizontal generator extends as a unit at
the toroidal boundary by the unipotent determinant calculation
above Lemma~\ref{ctwist:vertical}. On the whole model its divisor
is an integral multiple of the single reduced special fibre.

Every test extends to the toroidal model by properness. The
pullback of its connected semiabelian group is the given connected
semiabelian extension by uniqueness over a trait (SGA~7 I, IX,
Proposition 3.2). This is an isomorphism respecting the endomorphisms
by generic-fibre uniqueness. Hence its Lie eigensummands, including
those of a toric degeneration, are exactly the test lattices.
Pulling back the divisor in an integral local generator gives the
same normalized integer at every such test.

For general finite marking choose another full neat level of order
prime to \(\ell\) and a geometric connected component of the common
refinement lying over the original component. Over a finite field
extension it and both projection maps are defined. Fix one place
of this field above \(w\). The common component is finite étale
and surjective over each of its two characteristic-zero image
components, as follows by forgetting finite torsion bases. Its two
pulled-back determinant generators differ by a scalar in the
field of definition: both are horizontal for the same connection,
so their quotient has zero differential on the smooth geometrically
connected component. On the auxiliary prime-to-\(\ell\) component
the generator exists by determinant-one monodromy at its full neat
level, as proved above Lemma~\ref{ctwist:vertical}. Any test on the
original component first extends to the fixed base place and
then lifts generically to the common component after a finite
extension of its fraction field, by surjectivity. Connected
semiabelian pullbacks agree by the same uniqueness. The order on
the lifted test is the order on the auxiliary component plus the
valuation of the fixed scalar, both independent of the test after
normalization to \(F_w\). This proves the claim without requiring
an integral level-\(\ell\) marking of a semiabelian group.
\end{proof}

\begin{theorem}[Integral determinant comparison at the ramified primes]
\label{cmeig:splitmodel}
For the fields of Proposition~\ref{host:art}, under the same
classical moduli and test identifications as in
Proposition~\ref{cmeig:hyperspecial}, the common normalized order
exists at every finite place. For each indefinite component the
uncorrected arithmetic degree of the elementary horizontal
determinant is therefore \((p-q)A_\chi/2\).
\end{theorem}
\begin{proof}
Fix \(w,\tau\), and use a neat marking away from its residue
characteristic \(\ell\), with maximal level at \(\ell\) stabilizing
the single self-dual lattice (and its periodic translates).
As in Proposition~\ref{cmeig:hyperspecial} we can pass to this
level and to a finite local splitting field. At every place of
\(E^+\) above \(\ell\) the quadratic algebra is split or
unramified by Lemma~\ref{cmeig:row}. These are exactly the
unitary cases (AL) and (AU) of the Pappas--Rapoport splitting
model. Its order \(O_E\) is maximal, and its alternating lattice
has index one in its dual. These are the hypotheses at the
maximal level of \href{https://stephane-bijakowski.perso.math.cnrs.fr/PRgeometry.pdf}
{Bijakowski--Hernandez, \S2.4 and Theorem 2.30}, which gives a
smooth splitting model in these cases. At two the imaginary
quadratic subfield \(k\) already splits by Lemma~\ref{host:different},
so only the linear case (AL) occurs there.

Use the ordering of the \(E\)-embeddings in
\href{https://www.kwlan.org/articles/cpt-spl.pdf}
{Lan, Choices 2.3.1, Definition 2.3.3}: put \(\tau,\bar\tau\)
first as a consecutive pair in their \(E^+\)-place, and order the
other conjugate pairs consecutively. In the split case the pair
is in distinct integral field factors. In the inert unramified
quadratic case its restrictions to the maximal unramified subfield
are distinct, so the two are again separated by integral
idempotents over the splitting base. On these two different
summands the first quotient of the decreasing Hodge filtration
is labelled respectively by \(\tau,\bar\tau\).
In Lan's interleaved ordering the intervening graded piece on
the other summand is zero. In the increasing notation of
Bijakowski--Hernandez these are the two \emph{last} quotients.
Reversing the ordering within each list, with conjugate lists
paired term by term, is permitted in their \S2.4.

Here are the model comparisons used with the smoothness theorem.
For the single periodic self-dual lattice the open splitting model
parametrizes the polarized Hodge flags of the given ranks. Separating
the unramified idempotents turns Lan's interleaved flags into
the flags for the separate unramified embeddings; conversely they
are interleaved by direct sum. The polarization conditions can
also be compared on the common generic fibre, without identifying
their scheme structures on nonflat thickenings. Indeed both flag
problems are closed in the same relative product of flag varieties
over the naive moduli problem, by the endomorphism and pairing
conditions. They give the same unique flags on the generic fibre.
The Bijakowski--Hernandez model is smooth, hence equals the flat
closure of this generic fibre in that flag space, and is normal.
Lan's normalized open model is the normalization of that closure
on the selected generic components: the forgetful map of the
normalized underlying moduli to the naive moduli is finite and
generically open and closed at full lattice-stabilizer level
(Proposition 2.2.11). Finiteness persists after pullback to the
flag space and after normalization over the excellent base
(Definition 2.4.5 and Proposition 2.3.10). Consequently the normalized
open model on those components is the smooth one, by finite
birationality over a normal scheme.

Lan's Theorem 3.4.1 and Proposition 3.4.14 give a proper flat
toroidal compactification with geometrically regular fibres for
smooth cone data: geometric regularity of each fibre is
equivalent to that of the open splitting model. It is thus smooth.
By Proposition 3.1.2 and Definition 3.1.3 the extended filtration
is a filtration by subbundles on the Hodge bundle of the
semiabelian extension (identified with that of its dual by the
principal polarization, with conjugate labels if using the
dual-action convention). In particular the two locally free
first quotients above exist along the whole boundary.

For any trait write \(H\) for its integral invariant
differentials in the relevant unramified summand. If \(H'\) is
the kernel of its first quotient, then \(H'\) is saturated and
its generic fibre is the sum of all the other eigenspaces in
that summand. Hence
\[
   (H/H')^\vee=\{u\in H^\vee:u(H')=0\}
           =H^\vee\cap (H^\vee_{\rm frac})_\tau .
\]
The identical argument on the conjugate summand gives the
second test lattice. This uses the end quotient of the
filtration: in general an intermediate graded piece has a
different lattice. The determinant ratio of these dual quotients
is therefore an invertible sheaf which gives exactly the
ordinary test metric at \(\ell\). Its generic horizontal section
is a unit across the characteristic-zero boundary, as before.
Over the fixed geometric component the smooth proper model
has a single reduced special fibre. Lemma~\ref{ctwist:vertical}
and the common-refinement argument of
Proposition~\ref{cmeig:hyperspecial} give equality of all orders,
including on semiabelian tests. A different ordering or model
can be used for a different pair \(\tau,\bar\tau\), since orders
are calculated pairwise before multiplying sections.

Finally compare with the CM point of Proposition~\ref{cmeig:degree}
at each place and with its archimedean norm. Its finite extension
preserves normalized arithmetic degrees. This gives the claimed
global degree and hence the uncorrected product degree in
\eqref{ctwist:degree}.
\end{proof}

\paragraph{The selected good cover.}
The level condition on \(Y\) fixes its auxiliary marking to have
order one. Its finite marking is the half-marking at the odd integer
\(M\ge3\).
The characteristic-zero comparison can be made from the moduli
property. The universal object is proper and semiabelian, hence
abelian. Its rigidified line is relatively ample and its
polarization kernel is trivial, hence its polarization \(\lambda\)
is principal. The adjunction property for its biextension gives
the specified Rosati involution on \(O_E\). Conversely from these
PEL data pull the normalized Poincar\'e bundle back along
\((1,\lambda w_0)\). Its polarization is
\(\lambda(w_0+\bar w_0)=\lambda\); the rigidification, symmetry
and cube law are the canonical ones. Its biextension is the
one for \(\lambda\). This reconstruction also recovers the
original line from its graph condition. The rigidified isomorphism
is unique: an automorphism is multiplication by a unit on the
abelian scheme, which comes from the base by properness and
geometric connectedness and equals one at the identity.

Over the splitting field the determinant condition is precisely
the \(p,q\) multiplicity condition. On a reduced
characteristic-zero base the Hodge eigensummands are vector bundles
and their ranks can be tested at geometric points. The mark kills
\(\mathfrak m\) and on the complementary summand is the reduction
of a polarized Tate marking. By CRT these restrictions determine
it on the whole lattice. Its compatibility at all torsion orders
in the definition of an object is a condition on the local
lattice genus, not a choice of additional sections. More explicitly,
in a connected complex PEL family the Betti lattice, its pairing,
action and finite mark are locally constant; a compatible marked
Tate isometry at one point transports to every point. Equivalently
at a geometric point the reduction of a compatible tower is a
point of the prescribed compact-open marking orbit. Conversely
any representative of that orbit gives the required compatible
reductions at all \(n\). To check the roots of unity, fix the
primitive \(M\)-th root in the coefficient field and, at a geometric
point, extend it to a compatible system \((\zeta'_n)\). Relative
to the tower \((\zeta_n)\) of a Tate marking there is a unique
\(u\in\widehat{\mathbb Z}^{\times}\) with \(\zeta'_n=\zeta_n^u\)
for all \(n\ge1\). The reference lattice has an
integral similitude of multiplier \(u\) fixing the half-mark.
At \(\ell\mid M\), use \(1\) on the marked half and \(u_\ell\) on
its pairing-dual half. At \(\ell\nmid M\), use a central scalar of
norm \(u_\ell\) at each place of \(F_0\) over \(\ell\), using the
surjectivity of split and unramified unit norms in Lemma~\ref{cmeig:row}.
Precomposing a Tate marking by this similitude gives the prescribed
tower with the same marked sections. At order \(M\) the Weil pairing
with the marked basis and the fixed root then determines the basis
on the opposite half uniquely. This construction works over the
base at order \(M\) by perfect duality of finite étale groups. The
higher towers occur only in the geometric-point existence test.
This identifies the half-marking with the corresponding fine PEL
level on the selected components. After finite extension the
source moduli property gives a classifying map for the classical
universal object on its reduced smooth base, and the classical
moduli property gives the reverse map on the reduced smooth
source base. The reconstructed universal objects are isomorphic,
so both composites are the identity by uniqueness. The maps
preserve the underlying group and polarization, also after
restricting to a specified marked row.
\src{src/ReciprocityTheory/ReductiveGroups/ArithmeticGeometry/BirationalGeometry/ToroidalCompactification/TwoWeightEqOneAddRoot.lean}{131}{The compatible torsion isometries and primitive roots};
\src{src/ReciprocityTheory/ReductiveGroups/ArithmeticGeometry/BirationalGeometry/ToroidalCompactification/TwoWeightEqOneAddRoot.lean}{144}{The two marks, signature and geometric-point existence condition};
\src{src/ReciprocityTheory/ReductiveGroups/ArithmeticGeometry/BirationalGeometry/ToroidalCompactification/PushInvertible.lean}{29}{The universal object and its reduced-base moduli property};
\src{src/Development/AbelianVarietyTheory/Definitions/Divisors/NormNegCmConstantEqAtAllEmbeddings.lean}{187}{The selected level, period charts and polarization scale}.

On a test DVR the definition supplies a semiabelian group whose
generic fibre is cartesian over this universal group. Its fibres
are geometrically connected. The test vectors form a basis for
all integral eigenvectors at each labelled embedding, so they give
exactly the intersections used in Theorem~\ref{cmeig:splitmodel}.
Indeed the images of algebraic integers in the coefficient field
belong to the valuation ring; the equality expressing the
eigenvector property therefore holds integrally if and only if
it holds over the fraction field, by torsion-freeness of the
smooth Lie lattice. The theta generator in the test contributes
no factor to the parallel character. At infinity the comparison
uses the polarization metric \(\pi H\) on Lie vectors; this is
the metric used in Proposition~\ref{cmeig:degree}. Thus the
ordinary parallel determinant comparison follows on the selected
principal family with its given period readings.
\src{src/Formalism/LocalToGlobal/ProjectiveLattices/VAtVNeBot.lean}{53}{Cartesian semiabelian tests and the universal integral eigenbasis condition};
\src{src/ReciprocityTheory/ReductiveGroups/ArithmeticGeometry/SemiabelianSchemes/AbelianSchemeQuotients/IsAbelianScheme.lean}{41}{Connected geometric fibres of the semiabelian group};
\src{src/Development/AbelianVarietyTheory/Definitions/CMLattices/IsDiffOf.lean}{101}{Polarization norm in the CM comparison}.

\begin{corollary}[Short sections in the final weight on a good cover]
\label{cmeig:final}
Fix one good deep cover \(Y\) as in Corollary~\ref{crev:good}
on a selected host, with \(p,q\ge1\). Take the admissible final
weight \(\lambda\) and the level, denominator, field and bad-prime
size bounds of Theorem~\ref{bfam:oncharts}.
On each marked component \(Z\), its regular cusp rule space of
weight \(\lambda\) is spanned over the coefficient field by
integral sections \(s_i\) with
\[
     \log\|s_i\|_{\sigma}\le
          C\bigl((frL)^3 W+(frL)^5\bigr),\qquad
     W=\sum_{\tau,i}|\lambda_{\tau i}|,
\]
at every embedding in the theorem. The same constant works over
these choices under the common stage bounds.
\end{corollary}
\begin{proof}
Start with the integral generators of Corollary~\ref{crev:good}.
Multiply by the key of Theorem~\ref{ctwist:semi} and by the
product with exponents \(-j_\tau\) in \eqref{ctwist:degree}.
Theorem~\ref{cmeig:splitmodel} computes the latter degree at the
actual integral eigenlattices. Outside the excluded primes the
level is hyperspecial, so the common order is integral by
Proposition~\ref{cmeig:hyperspecial}. At the remaining primes
round the product order once as in \eqref{bfam:flat}.
Proposition~\ref{host:period} supplies the uniform numerical
degree bound. Both unit sections and their inverses preserve the
cusp condition by Corollary~\ref{crev:twist}. The final-weight
estimate of Theorem~\ref{bfam:oncharts} now gives the claim.
This argument starts with \(Y\) and the stated numerical bounds.
Selecting such covers compatibly at all stages and identifying the
restriction maps on the modified flags enter separately in the
Hilbert--Samuel estimate.
\end{proof}
